\appendices
\section{Map of Agent Communication Side-Channel Studies}\label{app:map}
\paragraph{Inclusion and retrieval}
A study is included if an LLM selects actions with external feedback able to affect later actions, and the study links that execution to a communication-side inference or an evaluated protection. Client, tool, model and inter-agent interfaces are eligible; transit, recipient and active observers remain separate. Whole semantic trajectories or browser instrumentation alone, deliberate covert messages and prompt-injection defenses without a communication-side endpoint are excluded~\cite{yao2026intentminer,uitraces2026}. The window covers first public submissions from 1 January 2024 through 9 November 2026. Twelve topic queries, a native arXiv query (180 records) and public-web queries (33 scholarly candidates) gave 209 records after deduplication; two citation leads brought screening to 211, of which fourteen were direct, forty adjacent, 155 outside scope and two unresolved. A retained published seed (CHeaT) gives fifteen direct sources with forty-four settings. Twelve of the fifteen lack a verified formal publication at the cutoff. Exact queries and dispositions are in Supplement S4.

\paragraph{Coding and reliability}
Each setting record binds source version and locator to access, features, learner, target, population, split, intervention, result and cost. We assigned and decided every code against the cited passage; AI tools located candidate passages and proposed drafts, which we accepted, changed or rejected. An independent security PhD researcher, blind to the manuscript and our codes, re-coded the decision-critical subset from the sources without AI tools. Access class (44/44), intervention purpose (9/9), and adaptation and utility in the four privacy settings (4/4 each) agree fully ($\kappa=1$); the protection flag agrees on 38/44 ($\kappa=.61$) and transfer relations on 5/6 ($\kappa=.77$). After adjudication of the seven disagreements against the sources, no privacy setting has both an adapted attacker and matched utility (Supplement S6.1).

\paragraph{The map}
Supplement S11 lists all forty-four settings with their interface, access, target, signal, unit, reported result and intervention evidence. Twenty-two use transit observations, ten recipient access, eleven active access and one inferred preferences. Three settings compare identity obfuscation and two adversarial-Agent blocking; thirty-four contain no protection comparison. Scores keep their source-defined metric and target and are not cross-paper rankings. Table~\ref{tab:transfer-matrix} relates the Agent evidence to established traffic-analysis conditions: a condition is \emph{retained} when it applies directly, \emph{extended} when an execution condition must be added, \emph{untested} when the Agent evidence does not evaluate it, and an \emph{access mismatch} when host instrumentation is mistaken for passive observation.

\begin{table*}[t]\centering\small
\caption{Transfer of established conditions to Agent execution. Relationship labels concern the stated mapping; evidence strength and missing tests appear separately. A missing Agent test does not refute the classical result.}\label{tab:transfer-matrix}
\setlength{\tabcolsep}{4pt}
\begin{tabularx}{\textwidth}{@{}P{.25\textwidth}P{.29\textwidth}P{.10\textwidth}Y@{}}\toprule
Established relation and condition & Agent mapping / consequence & Relation & Agent evidence and remaining test \\\midrule
Unknown boundaries and background affect event recognition~\cite{wang2016realistic,guan2021bapm,deng2023ares,trimananda2020pingpong}. & Tool bursts, model waits and task intervals are different units. Evaluate proposed events with false alarms as well as classification. & Extended & ANT supplies stage boundaries; the concurrent feedback study (Supplement S10) tests particular detectors. Greater Agent localization difficulty is untested. \\\addlinespace[3pt]
Defense evaluation must allow adaptation to the protected distribution~\cite{dyer2012peekaboo,mathews2023critical,tramer2020adaptive}. & Record protected training access and still-visible links; an unchanged learner family can refit. & Retained & Transformed and live OHTTP settings report adaptation in one preprint~\cite{dangol2026workflow}; WRA's matched protected examples remain unresolved. \\\addlinespace[3pt]
Coarse timing/volume can remain useful under specified padding~\cite{dyer2012peekaboo}; serving streams expose length/gap signals~\cite{weiss2024prompt,whisperleak2025}. & Name the shaping policy, receiving hop and Agent workload before carrying over a measured effect. & Untested generalization & Cloud and ANT measure model-view recognition~\cite{pouryousef2026traffic,li2026ant}. A general protected-Agent SSE claim is not tested by these observations. \\\addlinespace[3pt]
An observation result is tied to the visible paths and representations~\cite{chen2010web,apthorpe2019stp}. & Feedback forwards tool information into later model requests across separately controlled services; evaluate every permitted residual view. & Extended & Published WRA tool-side evidence and preprint model/recipient results establish different exposed interfaces. Causal feedback attribution needs its own test. \\\addlinespace[3pt]
STP tunnel endpoints remove cover before it reaches devices or servers~\cite[Sec.~7.1]{apthorpe2019stp}. & When cover executes as research, add original-task resource, state and feedback constraints to the traffic budget. & Extended & Published WRA reports executed research decoys and utility under a different intervention~\cite{jeong2026wra}. Our controlled LLM case separates task completion from state integrity (Section~\ref{sec:state}); it does not evaluate WRA writes. \\\addlinespace[3pt]
Host traces can reconstruct distributed request paths~\cite{reynolds2006wap5,tak2009vpath}. & Instrumentation can provide ground truth; a transit attacker must infer from its own permitted observations. & Access mismatch & ANT stage labels and local task intervals are evaluator-provided. Host tracing alone cannot establish passive encrypted-stream localization. \\\bottomrule
\end{tabularx}
\end{table*}



\section{Defense Design Space}\label{app:defense}
Table~\ref{tab:defense-space} lists defense layers, what each requires the defender to control, the costs to measure together and the exposure that remains.

\begin{table*}[t]\centering\small
\caption{Defense design space: enforcement authority, costs to measure and remaining exposure. Source-specific measurements appear in the text and Supplement S7.}
\label{tab:defense-space}
\setlength{\tabcolsep}{4pt}
\begin{tabularx}{\textwidth}{@{}P{.17\textwidth}P{.25\textwidth}P{.23\textwidth}Y@{}}\toprule
Layer / mechanism & What must be trusted or controlled & Costs to measure together & Residual exposure and practical decision \\\midrule
Naming / association: encrypted name resolution, relays, OHTTP & Name resolver or relay placement; the parties able to link the user to the destination & Added hops, bytes, latency and availability & Shape can survive name hiding. Choose this layer when destination association is the private target; then test remaining observers. \\\addlinespace[3pt]
Representation: message, TLS-record or packet padding & Component that controls the final visible size; padding survives later encoding & Added bytes, record/packet count, fragmentation & Timing, count and other links remain. State the padding layer and test downstream captures. \\\addlinespace[3pt]
Scheduling: batching, pacing, fixed-rate cover & Sender/relay controls output schedule and cover lifetime & First-response delay, completion delay, bandwidth during idle periods & Admission, termination and overflow can reveal activity. Budget these events as part of the visible trace. \\\addlinespace[3pt]
Transport cover: dummy exchanges discarded at a tunnel endpoint & Trusted cover terminator; no dispatch into business tools & Cover bandwidth, queueing, resource use & Other endpoints still see their inputs. Prefer this boundary when the aim is traffic shaping without extra business work. \\\addlinespace[3pt]
Execution cover: extra research or tool calls & Tool privileges, isolated state and whether decoy feedback reaches the planner & Model tokens, tool quota, task success, state changes and latency & A decoy can affect the real execution. Measure privacy and original-task utility on the same protected runs. \\\addlinespace[3pt]
Shared-state isolation / minimization & Tenant separation, cache keys and information passed between services & Lost cache reuse, compute, latency, task accuracy & Isolation removes a sharing channel, not every traffic signal. Minimize tool arguments and feedback according to each recipient's role. \\\bottomrule
\end{tabularx}
\end{table*}

